% Sección I: Introducción, objetivos y correspondencia normativa
% evidencia:E01,E02,E03

\section{Introducción}

La gestión operativa de un hospital de alta complejidad exige articular la anticipación de llegadas diarias con el comportamiento de la permanencia intrahospitalaria y la distribución horaria de los recursos. En el Hospital Universidad del Norte (HUN), la disponibilidad de un repositorio estructurado bajo el estándar Observational Medical Outcomes Partnership Common Data Model (OMOP CDM) v5.4~\cite{ohdsi54} permite analizar trayectorias agregadas sin exponer registros individuales de pacientes; sin embargo, dicha representación administrativa carece de tablas de capacidad física instalada, camas censales habilitadas y dotación contractual por turno~\cite{mediflow2026informe}. MediFlow se concibe ante esa restricción como una herramienta de apoyo a la planificación bajo el hilo conductor \textbf{comprender, predecir y decidir}: reconstruye el flujo histórico por etapas, evalúa la incertidumbre de la demanda con validación temporal estricta y traduce pronósticos y reglas de permanencia en escenarios paramétricos interpretables con supuestos explícitos.

El aporte defendible del proyecto reside en la integración trazable de cuatro capas desacopladas de los datos individuales en despliegue: extracción SQL agregada de solo lectura en origen, modelado predictivo con validación de origen móvil y explicabilidad aditiva, análisis operacional basado en colas y balance de flujo, e interfaces interactivas (aplicación web FastAPI/React y tablero histórico en Power BI). El sistema no afirma haber demostrado una reducción clínica observada de tiempos de espera ni una asignación óptima sobre recursos reales del hospital; todo escenario cuantificado se presenta como hipótesis operacional sujeta a contraste con capacidad física y marcas temporales del sistema de información hospitalario (HIS).

\subsection{Objetivos y alcance editorial}

El objetivo general consiste en analizar el flujo de pacientes y apoyar la planificación operacional del HUN mediante pronósticos de demanda con intervalos de incertidumbre y simulaciones interpretables bajo supuestos explícitos. Se desprenden cinco objetivos específicos: (a) reconstruir la demanda por servicio, las etapas de atención, las transiciones entre áreas y el censo horario desde agregados OMOP; (b) auditar la calidad, la completitud, la consistencia temporal y el aporte empírico de factores externos mediante protocolos anti-fuga; (c) contrastar modelos predictivos de corto y largo horizonte con métricas libres de explosión por ceros y evaluación temporal estricta; (d) cuantificar escenarios paramétricos de dotación equivalente, permanencia promedio y redistribución intradía de egresos; y (e) comunicar las capacidades implementadas, las diferencias entre versiones experimentales y los límites para una eventual validación clínica.

\subsection{Correspondencia con los requisitos oficiales del certamen}

Las exigencias normativas provienen exclusivamente de los términos y condiciones de participación y del enunciado técnico oficial del Data Challenge HUN 2026~\cite{hun2026terminos,hun2026reto}. Los entregables oficiales comprenden la documentación técnica (metodología, análisis, resultados, conclusiones y recomendaciones), un repositorio en GitHub con el código autorizado, un enlace interactivo en Power BI y la sustentación oral cuando corresponda. El enunciado permite abordar uno o varios componentes del flujo y no exige emplear todas las tablas clínicas como predictores ni desplegar un piloto asistencial previo.

\begin{table}[htbp]
\caption{Criterios oficiales de evaluación y desempates del Data Challenge HUN 2026}
\label{tab:criterios_oficiales}
\centering
\scriptsize
\begin{tabularx}{\columnwidth}{@{}L{0.44\columnwidth} r Y@{}}
\toprule
Criterio oficial de evaluación & Peso & Bloque principal del informe \\
\midrule
Comprensión del problema y flujo hospitalario & 15\,\% & Sec. I (Intro.); Sec. II (Datos y flujo) \\
Calidad, preparación e integración de datos & 15\,\% & Sec. II--III (Flujo, calidad y externos) \\
Modelo analítico/predictivo y rigor metodológico & 20\,\% & Sec. IV--V (Método y resultados) \\
Optimización y propuestas de mejora operacional & 20\,\% & Sec. VI (Optimización paramétrica) \\
Integración, visualización y dashboard & 15\,\% & Sec. VII (API, web y Power BI) \\
Ingeniería de software, reproducibilidad y doc. & 5\,\% & Sec. VIII (Ingeniería y privacidad) \\
Presentación, interpretación y recomendaciones & 10\,\% & Sec. IX--X (Discusión y conclusiones) \\
\midrule
Total de la rúbrica oficial & 100\,\% & Cobertura transversal y Anexos A--D \\
\bottomrule
\end{tabularx}
\ieeenota{Basado en~\cite{hun2026terminos} (p.~4, secs. 13--14) y~\cite{hun2026reto} (p.~3). Según la sec.~16 de~\cite{hun2026terminos} (p.~5), los empates se resuelven sucesivamente por optimización y mejora operacional, modelo analítico y rigor metodológico, comprensión del flujo hospitalario y valoración conjunta de utilidad, factibilidad y carácter accionable. La maquetación en \LaTeX\ con clase \texttt{IEEEtran} es una decisión editorial de los autores y no un requisito oficial.}
\end{table}

La Tabla~\ref{tab:criterios_oficiales} sintetiza los pesos de la rúbrica oficial y su correspondencia orientadora con las secciones de este informe, manteniendo explícita desde el inicio la separación entre corridas experimentales archivadas, metadatos inspeccionados del artefacto web y verificaciones locales.
